\section{重心（CoM）に基づく歩行制御の考え方}

重心（CoM）は，ロボット全体の質量分布を代表する点として定義される．二足歩行ロボットにおいては，多数のリンクと関節から構成される全身の運動を直接扱うことは複雑であるため，CoMの運動に着目して歩行の安定性を議論する手法が広く用いられている．

歩行中のロボットの安定性は，CoMの位置や加速度が床反力とどのような関係にあるかによって特徴づけられる．特に，CoMの水平方向の運動は，転倒方向や支持脚への負荷分布に大きく影響する．このため，多くの歩行制御手法では，CoMの目標軌道を生成し，それに追従させることが基本的な方針となっている．

一方で，CoMの鉛直方向の運動は，水平方向に比べて制御対象として明示的に扱われない場合が多い．これは，従来のモデルではCoM高さを一定と仮定することで，問題を簡略化してきたためである．しかし，実環境においては，床の変形や沈み込みによって，CoM高さが変動する状況が生じる．この点については，第4章以降で詳しく議論する．